\section*{Source cofinal realization and terminal extinction}

The realization governing this chapter is the spectral--polar complement of the common enriched antecedent. Let
\[
 \mathcal D_R=C_c^\infty((-R/2,R/2)),\qquad
 \mathcal D_R^{\rm cof}:=\mathcal D_R,\qquad
 \widehat\psi(t)=\int_{\mathbb R}\psi(x)e^{-itx}\,dx ,
\]
and, for \(\psi,\phi\in\mathcal D_R\),
\[
 h_{\psi,\phi}(z)
 =\widehat\psi(z)\overline{\widehat\phi(\overline z)},\qquad
 \check h_{\psi,\phi}(a)
 =\langle\psi,T_a\phi\rangle ,
 \qquad (T_a\phi)(x)=\phi(x-a).
\]

Let \(P_+^{\rm sh},P_-^{\rm sh}\) be the orthogonal projectors of the two sheets, and fix
\begin{equation}\label{eq:amp-rh-owner-q}
 q=\frac59,\qquad
 A=P_+^{\rm sh}+qP_-^{\rm sh},\qquad
 D=\sqrt{1-q^2}\,P_-^{\rm sh},\qquad A^*A+D^*D=I.
\end{equation}
If \(\Xi_R\) is the restriction of the common antecedent, \(Q_R=I-P_R\) its orthogonal complement, and \(\jmath_-\) the isometric inclusion in the antispecular sheet, construct
\begin{equation}\label{eq:amp-rh-owner-binding}
 E_R:=\jmath_-Q_R\Xi_R,\qquad
 AE_R=qE_R,\qquad
 E_R^*E_R=\mathcal B_{{\rm W},R}.
\end{equation}
The conservation identity is iterated without deleting the last summand:
\begin{equation}\label{eq:amp-rh-owner-telescopia}
 \mathcal B_{{\rm W},R}
 =\sum_{n=0}^{N-1}(DA^nE_R)^*(DA^nE_R)
 +(A^NE_R)^*(A^NE_R).
\end{equation}
The terminal is computed before taking the limit:
\begin{equation}\label{eq:amp-rh-owner-terminal}
 (A^NE_R)^*(A^NE_R)
 =q^{2N}\mathcal B_{{\rm W},R}\longrightarrow0
 \qquad(N\longrightarrow\infty).
\end{equation}
Only after this extinction do we define
\begin{equation}\label{eq:amp-rh-owner-l}
 \mathscr L_Rf:=(DA^nE_Rf)_{n\geq0},
 \qquad
 \boxed{\mathcal B_{{\rm W},R}=\mathscr L_R^*\mathscr L_R\succeq0}.
\end{equation}
Restriction compatibility publishes these factorizations as a single cofinal form:
\begin{equation}\label{eq:amp-rh-owner-publicacion}
 B_W=\Xi^*(I-P)\Xi
 =\bigl((I-P)\Xi\bigr)^*\bigl((I-P)\Xi\bigr)\succeq0,
 \qquad
 B_W(f,f)=\|(I-P)\Xi f\|^2.
\end{equation}

\begin{theorem}[Cofinal Riemann--Weil closure]
\label{thm:amp-rh-cierre-cofinal-propietario}
The complete Guinand--Weil form is positive on the cofinal class \(\bigcup_{R>0}\mathcal D_R^{\rm cof}=C_c^\infty(\mathbb R)\). By Weil's reversible criterion, every nontrivial zero \(\rho\) of the Riemann zeta function satisfies
\[
 \Re\rho=\frac12.
\]
\end{theorem}

\begin{proof}
The constraints of \eqref{eq:amp-rh-owner-binding} turn the terminal of \eqref{eq:amp-rh-owner-telescopia} exactly into the left-hand side of \eqref{eq:amp-rh-owner-terminal}. Its geometric extinction permits passage to \eqref{eq:amp-rh-owner-l}; cofinal naturality produces \eqref{eq:amp-rh-owner-publicacion}. The prime--gamma--polar identification with the complete functional and the cofinal separation from the principal source \eqref{eq:realizaciones-rh-identificacion-gw} allow Weil's criterion to be applied. Subsequent calculations check that publication again in an auxiliary regulated chart, without becoming premises of the theorem.
\end{proof}

\section*{Auxiliary common-regulator control and three-channel identity}

The following regulated calculation expands the Gram difference from its constituent operators before assessing its sign. Its residual is a \texttt{CONTROL\_NON\_GOVERNING}: it delimits this particular window and neither replaces nor conditions the cofinal realization \eqref{eq:amp-rh-owner-binding}--\eqref{eq:amp-rh-owner-publicacion}. Fix a function \(\vartheta\in C_c^\infty([0,\infty))\) such that \(0\leq\vartheta\leq1\), \(\vartheta=1\) on \([0,1]\), and \(\vartheta=0\) on \([2,\infty)\), and set
\[
 \vartheta_L(a)=\vartheta(a/L),\qquad \vartheta_L(0)=1 .
\]
It is the only length regulator: it applies simultaneously to the prime-power clocks and the gamma channel; the scalar and polar terms are its zero-length atoms. Regularizing each sheet with a different cutoff is not allowed.

For
\[
 \ell_{p,k}=k\log p,\qquad
 w_{p,k}=(\log p)p^{-k/2},\qquad
 \mu_\Gamma(a)=\frac{e^{-a/2}}{1-e^{-2a}},
\]
let us also define
\begin{align*}
 c_\Gamma&=\psi_\Gamma(1/4)-\log\pi<0,\\
 A_\psi&=\widehat\psi(i/2),&
 B_\psi&=\widehat\psi(-i/2),&
 b_\pm(\psi)&=\frac{A_\psi\pm B_\psi}{\sqrt2}.
\end{align*}
At level \(N_m=9^m\), let \(j_m\) be the uniform inclusion and \(\widehat\delta_{a,m}\) the normalized nonadic carry channel. Their moments are computed directly from the odometer:
\begin{equation}\label{eq:amp-rh-momentos-odometro-explicitos}
 j_m^*j_m=I,\qquad
 \widehat\delta_{a,m}^*\widehat\delta_{b,m}
 =(I-T_a)^*(I-T_b).
\end{equation}
In particular, \(\widehat\delta_{a,m}\) depends on the seam and on \(T_a\), not on the zeta function or its zeros.

The two regulated columns are
\begin{align}
 \mathcal C^+_{m,R,L}\psi
 ={}&
 \Bigl(
  (\sqrt{\vartheta_L(\ell_{p,k})w_{p,k}}\,
       \widehat\delta_{\ell_{p,k},m}\psi)_{p,k},
 \nonumber\\[-2mm]
 &\hspace{11mm}
  a\longmapsto
  \sqrt{\vartheta_L(a)\mu_\Gamma(a)}\,
       \widehat\delta_{a,m}\psi,\,
  \zeta_+b_+(\psi)
 \Bigr),                                             \label{eq:amp-rh-columna-mas-explicita}\\
 \mathcal C^-_{m,R,L}\psi
 ={}&
 \Bigl(
  (\sqrt{2\vartheta_L(\ell_{p,k})w_{p,k}}\,j_m\psi)_{p,k},\,
  \sqrt{-c_\Gamma}\,j_m\psi,\,
  \zeta_-b_-(\psi)
 \Bigr),                                             \label{eq:amp-rh-columna-menos-explicita}
\end{align}
where the sums run only over lengths for which \(\vartheta_L(\ell_{p,k})\neq0\), and \(\zeta_\pm\) are the orthogonal polar lines of the cyclic representation of order \(12\).

\subsection*{Indecomposable coupling prior to the sign}

The relation between \eqref{eq:amp-rh-columna-mas-explicita} and \eqref{eq:amp-rh-columna-menos-explicita} is not obtained by postulating a contraction between two Gram matrices already computed. Let
\[
 P_m=\iota_mE_m,\qquad Q_m=I-P_m
\]
be the explicit nonadic expectation. For the unitary clock \(\mathscr U_{\ell,m}\), write
\[
 a_{\ell,m}=P_m\mathscr U_{\ell,m}P_m,\qquad
 c_{\ell,m}=Q_m\mathscr U_{\ell,m}P_m ,
\]
so that the block unity gives
\begin{equation}\label{eq:amp-rh-defecto-schwarz-reloj}
 P_m-a_{\ell,m}^*a_{\ell,m}
 =c_{\ell,m}^*c_{\ell,m}\succeq0.
\end{equation}
For \(r=p^{-1/2}\), set
\[
 s_{p,m}=\frac{N_mr}{1+(N_m-1)r},
 \qquad \xi_{p,m}=\operatorname{arsech}(s_{p,m}),
\]
and let \(\mathscr S_{\xi_{p,m}}\) be the unitary polar colligation obtained by the Potapov--Ginzburg transform. Its Redheffer composition with the clock satisfies
\begin{align}
 C_{p,m}&=(s_{p,m}I-a_{\log p,m})
          (I-s_{p,m}a_{\log p,m})^{-1},\nonumber\\
 I-C_{p,m}^*C_{p,m}
 &=\bigl(1-s_{p,m}^2\bigr)
   (I-s_{p,m}a_{\log p,m}^*)^{-1}\nonumber\\
 &\quad{}\times c_{\log p,m}^*c_{\log p,m}
   (I-s_{p,m}a_{\log p,m})^{-1}.
   \label{eq:amp-rh-redheffer-defecto}
\end{align}
The Neumann expansion, before any spectral publication, gives
\begin{equation}\label{eq:amp-rh-torre-prima-explicita}
 \lambda_{p,m}(I-C_{p,m}^*C_{p,m})
 =\sum_{k\geq1}(\log p)p^{-k/2}
  \bigl(2I-T_{k\log p}-T_{k\log p}^*\bigr),
\end{equation}
with
\[
 \lambda_{p,m}
 =\frac{(\log p)r(1+r)}
 {(1-r)\kappa_{p,m}},
 \qquad
 \kappa_{p,m}
 =\frac{(1-s_{p,m}^2)(N_m-1)}
 {[N_m-(N_m-1)s_{p,m}]^2}.
\]
Thus the powers of a prime are coefficients of a single modular feedback, not a list selected afterwards. Let \(D_{p,L}\) be the diagonal contraction of this feedback's reachability port, multiplying its coordinate \(k\) by \(\sqrt{\vartheta_L(k\log p)}\). Regulation alters neither the clock nor its coefficients: it simultaneously compresses the output coordinates of the same tower.

The finite role port is not diagonal either. In the active sectors
\[
 \mathcal H_+
 =\operatorname{span}\{f_3,f_6,f_9\},\qquad
 \mathcal H_-=\operatorname{span}\{f_4,f_8\},
\]
the coisometry
\begin{equation}\label{eq:amp-rh-acoplamiento-angular}
 C_\angle
 =|f_4\rangle\langle f_3|
  +|f_8\rangle\langle f_9|
\end{equation}
satisfies
\[
 C_\angle C_\angle^*=I_{\mathcal H_-},\qquad
 I_{\mathcal H_+}-C_\angle^*C_\angle=P_6.
\]
The isometric chart \(U_{W_{24}\to12}\), obtained from the \(5\) incident octads and their Gram matrix \(4I_5+\frac43J_5\), transports this coisometry to the HMT boundary:
\[
 C_{W_{24}}
 =U_{W_{24}\to12}^*C_\angle U_{W_{24}\to12}.
\]

\begin{theorem}[Control of the Paley--Wiener--polar network and its port residual]
\label{thm:amp-rh-acoplamiento-prospectivo}
For fixed \(m,R,L\), define
\begin{align}
 \mathscr R^{\rm pr}_{m,L}
 &=
 \mathop{\star}_{\log p\leq2L}
 D_{p,L}\bigl(\mathscr S_{\xi_{p,m}}\star
              \mathscr U_{\log p,m}\bigr),\nonumber\\
 \mathscr R^\Gamma_{m,L}
 &=
 \int_0^{2L}{}^\oplus
 \mathscr U_{a,m}
 \vartheta_L(a)\,d\mu_\Gamma(a),\nonumber\\
 \mathscr C^{\rm TPK}_{m,R,L}
 &=C_{W_{24}}\star
 \bigl(\mathscr R^{\rm pr}_{m,L}
       \oplus\mathscr R^\Gamma_{m,L}\bigr).
 \label{eq:amp-rh-red-prospectiva}
\end{align}
This network is passive. In its ports' native energy spaces, let \(J_{+,m,R,L}:=\mathcal C^+_{m,R,L}\), \(J_{-,m,R,L}:=\mathcal C^-_{m,R,L}\), and \(\mathscr C^{\rm src}_{m,R,L}\) be the external contractive transfer of \eqref{eq:amp-rh-red-prospectiva}. Define
\begin{equation}\label{eq:amp-rh-publicaciones-construidas}
 \begin{aligned}
 \mathcal E_{m,R,L}
 &:=\mathscr C^{\rm src}_{m,R,L}J_{+,m,R,L}
    -J_{-,m,R,L},\\
 \mathscr D_{m,R,L}
 &:=\bigl(I-(\mathscr C^{\rm src}_{m,R,L})^*
                  \mathscr C^{\rm src}_{m,R,L}\bigr)^{1/2},\\
 \mathsf L_{m,R,L}&:=\mathscr D_{m,R,L}J_{+,m,R,L}.
 \end{aligned}
\end{equation}
Residual \(\mathcal E_{m,R,L}\) measures whether this auxiliary window chart alone reproduces the publication already constructed by the cofinal source. Its vanishing is neither a requirement of Riemann--Weil closure nor a premise of its positive factorization. No coefficient of \eqref{eq:amp-rh-red-prospectiva} depends on a zero of \(\zeta\), on Weil positivity, or on the sign to be obtained.
\end{theorem}

\begin{proof}
Every clock \(\mathscr U_{a,m}\) is unitary. Identity \eqref{eq:amp-rh-defecto-schwarz-reloj} identifies its carry port. The polar transform is unitary, and \eqref{eq:amp-rh-redheffer-defecto} proves passivity after the internal port is closed. The Redheffer composition of finitely many passive colligations is again passive. The gamma channel is obtained as a limit of direct sums in \([\varepsilon,2L]\); the estimate \(\mu_\Gamma(a)\|(I-T_a)\psi\|^2=O(a)\) when \(a\downarrow0\) gives the limit in the form domain.

Expansion \eqref{eq:amp-rh-torre-prima-explicita} produces exactly the prime-power coefficients. The direct integral produces the second entry of \eqref{eq:amp-rh-columna-mas-explicita}. The Hadamard transform at the polar port produces \(b_+\) and \(b_-\), and \eqref{eq:amp-rh-acoplamiento-angular}, conjugated by \(U_{W_{24}\to12}\), fixes the \(3\) inputs and the \(2\) outputs of the same state. These operations construct the transfer and the two native columns; their output difference is exactly \(\mathcal E_{m,R,L}\) in \eqref{eq:amp-rh-publicaciones-construidas}. Passivity guarantees that \(\mathscr D_{m,R,L}\) is well defined, but does not by itself cancel that auxiliary residual; nor does it need to do so for the cofinal closure already proved.

This network is not \(C_\angle\otimes I\): factors \((I-s_{p,m}a_{\log p,m})^{-1}\) mix each polar port with all powers of the clock, and the direct integral mixes the same port with the gamma continuum. Removing those terms changes the reachability column and is not a simplification of \eqref{eq:amp-rh-red-prospectiva}.
\end{proof}

\subsection*{Term-by-term expansion of the Gram difference}

\begin{lemma}[Prime--gamma--polar identity with a common regulator]
\label{lem:amp-rh-expansion-guinand-weil}
For \(\psi,\phi\in\mathcal D_R\),
\begin{equation}\label{eq:amp-rh-diferencia-columnas-regulada}
 \langle J_{+,m,R,L}\psi,J_{+,m,R,L}\phi\rangle
 -\langle J_{-,m,R,L}\psi,J_{-,m,R,L}\phi\rangle
 =
 \mathscr I_{{\rm GW},L}(\psi,\phi),
\end{equation}
where
\begin{align}
 \mathscr I_{{\rm GW},L}(\psi,\phi)
 ={}&
 h_{\psi,\phi}(i/2)+h_{\psi,\phi}(-i/2)
 +c_\Gamma\check h_{\psi,\phi}(0)                 \nonumber\\
 &+\int_0^\infty\vartheta_L(a)\mu_\Gamma(a)
 \bigl(
 2\check h_{\psi,\phi}(0)
 -\check h_{\psi,\phi}(a)
 -\check h_{\psi,\phi}(-a)
 \bigr)\,da                                       \nonumber\\
 &-\sum_{p,k}\vartheta_L(\ell_{p,k})w_{p,k}
 \bigl(
 \check h_{\psi,\phi}(\ell_{p,k})
 +\check h_{\psi,\phi}(-\ell_{p,k})
 \bigr).                              \label{eq:amp-rh-gw-truncado-explicito}
\end{align}
This is the complete truncated Guinand--Weil functional: it contains the polar term, the gamma factor, and the \(2\) orientations of every prime power under a single regulator.
\end{lemma}

\begin{proof}
By definitions \eqref{eq:amp-rh-columna-mas-explicita}--\eqref{eq:amp-rh-columna-menos-explicita},
\begin{align}
 \langle J_+\psi,J_+\phi\rangle
 -\langle J_-\psi,J_-\phi\rangle
 &=
 \langle\mathcal C^+\psi,\mathcal C^+\phi\rangle
 -\langle\mathcal C^-\psi,\mathcal C^-\phi\rangle .
 \label{eq:amp-rh-expansion-inicial}
\end{align}
In the prime channel, \eqref{eq:amp-rh-momentos-odometro-explicitos} gives
\begin{align*}
 &\langle(I-T_\ell)\psi,(I-T_\ell)\phi\rangle
 -2\langle\psi,\phi\rangle\\
 &\hspace{18mm}
 =-\check h_{\psi,\phi}(\ell)
  -\check h_{\psi,\phi}(-\ell).
\end{align*}
Multiplying by \(\vartheta_L(\ell_{p,k})w_{p,k}\) and summing produces exactly the final line of \eqref{eq:amp-rh-gw-truncado-explicito}.

In the gamma channel,
\[
 \langle(I-T_a)\psi,(I-T_a)\phi\rangle
 =2\check h(0)-\check h(a)-\check h(-a).
\]
Moreover, the integral representation of the logarithmic derivative of gamma gives, for \(t\in\mathbb R\),
\begin{equation}\label{eq:amp-rh-digamma-longitudes}
 \Re\psi_\Gamma\!\left(\frac14+\frac{it}{2}\right)-\psi_\Gamma(1/4)
 =2\int_0^\infty\mu_\Gamma(a)(1-\cos at)\,da .
\end{equation}
Integrating against \(h(t)/(2\pi)\) and using Fourier inversion yields
\begin{align*}
 &-\log\pi\,\check h(0)
 +\frac1{2\pi}\int_{\mathbb R}h(t)
  \Re\psi_\Gamma\!\left(\frac14+\frac{it}{2}\right)\,dt\\
 &\qquad =
 c_\Gamma\check h(0)
 +\int_0^\infty\mu_\Gamma(a)
  [2\check h(0)-\check h(a)-\check h(-a)]\,da .
\end{align*}
The regulated version applies the common substitution to this identity
\[
 \mu_\Gamma(a)\longmapsto
 \vartheta_L(a)\mu_\Gamma(a).
\]

Finally,
\begin{align*}
 \langle b_+(\psi),b_+(\phi)\rangle
 -\langle b_-(\psi),b_-(\phi)\rangle
 &=
 A_\psi\overline{B_\phi}
 +B_\psi\overline{A_\phi}\\
 &=h_{\psi,\phi}(i/2)+h_{\psi,\phi}(-i/2).
\end{align*}
The sum of the \(3\) expansions is \eqref{eq:amp-rh-gw-truncado-explicito}. Every equality has been computed before knowing the sign of their sum.
\end{proof}

\section*{Convergence, cofinality, and separation}

\begin{proposition}[Removing the regulator and cofinal naturality]
\label{prop:amp-rh-convergencia-cofinal}
For every \(\psi,\phi\in\mathcal D_R\), the limit
\[
 \lim_{L\to\infty}\mathscr I_{{\rm GW},L}(\psi,\phi)
 =\mathscr I_{\rm GW}(\psi,\phi)
\]
exists and coincides with the complete Guinand--Weil form. If \(R\leq R'\), inclusions \(\mathcal D_R\hookrightarrow\mathcal D_{R'}\) and refinements \(m\mapsto m+1\) separately intertwine \(J_{+,m,R,L}\) and \(J_{-,m,R,L}\). Thus identities \eqref{eq:amp-rh-diferencia-columnas-regulada} form a single compatible family, not a collection of window-dependent factorizations.
\end{proposition}

\begin{proof}
Since \(\check h_{\psi,\phi}\in C_c^\infty((-R,R))\), the prime-power sum remaining after cancellation of the scalar counterterms is locally finite: every term with \(\ell_{p,k}\geq R\) vanishes. In the gamma integral,
\[
 2\check h(0)-\check h(a)-\check h(-a)=O(a^2)
 \quad(a\downarrow0),
\]
while \(\mu_\Gamma(a)\sim(2a)^{-1}\); the integrand is \(O(a)\). For \(a>R\), both translations disappear and \(\mu_\Gamma(a)=O(e^{-a/2})\). Dominated convergence allows \(L\to\infty\).

Moment identity \eqref{eq:amp-rh-momentos-odometro-explicitos} does not depend on \(m\); it therefore defines the refinement isometry \(\widehat\delta_{a,m}\mapsto\widehat\delta_{a,m+1}\), while \(j_m\mapsto j_{m+1}\). The polar lines and \(U_{W_{24}\to12}\) remain fixed. As the window expands, the new clocks have disjoint translations on \(\mathcal D_R\) and add the same term \(2w_{p,k}\langle\psi,\phi\rangle\) to both columns before cancellation. This proves both intertwining relations.
\end{proof}

\paragraph{Finite control inherited from the nonadic fibre.}
Let \(\mathscr F_{729}\) be the finite local factor with \(729\) states, \(P^{[729]}\) the projector onto the \(81\) coarse states, and \(Q^{[729]}=I_{\mathscr F_{729}}-P^{[729]}\). Then
\begin{equation}\label{eq:amp-rh-729-81-648}
 729=81+648,\qquad
 \operatorname{rank}P^{[729]}=81,\qquad
 \operatorname{rank}Q^{[729]}=648.
\end{equation}
This census precedes the sign criterion and does not identify complement \(Q^{[729]}\) with connector residual \(\mathcal E_{m,R,L}\).

\begin{proposition}[Cofinal compatibility of the finite complement]
\label{prop:amp-rh-compatibilidad-complemento-finito}
\label{prop:amp-rh-residual}
The common-publication identities of \eqref{eq:realizaciones-rh-publicaciones-comunes}--\eqref{eq:realizaciones-rh-bw} form a single compatible family under refinements \(m\mapsto m+1\) and inclusions \(\mathcal D_R\hookrightarrow\mathcal D_{R'}\). In particular, the difference of the two Gram matrices is transported as the form of complement \((I-P)\Xi\), and census \eqref{eq:amp-rh-729-81-648} remains a finite control of the fibre.
\end{proposition}

\begin{proof}
The orthogonality of \(P\) and the common publications give \eqref{eq:realizaciones-rh-bw}. The \cref{prop:amp-rh-convergencia-cofinal} separately intertwines both columns under refinement and inclusion; subtracting them preserves the complement identity throughout the cofinal family.
\end{proof}

This auxiliary chart does not assert \(\mathcal E_{m,R,L}=0\): the \cref{prop:amp-rh-obstruccion-ventana-corta} delimits only that regulated model. The recovered census belongs to the source finite complement, not to the short-window control; hence the residual neither reopens nor conditions the cofinal theorem.

\begin{lemma}[The cofinal class separates spectral evaluations]
\label{lem:amp-rh-separacion-cofinal}
We have equality, not merely density,
\[
 \bigcup_{R>0}\mathcal D_R=C_c^\infty(\mathbb R).
\]
Moreover, for any distinct complex points \(z_1,\ldots,z_n\), the map
\[
 \mathcal D_R\longrightarrow\mathbb C^n,\qquad
 \psi\longmapsto
 (\widehat\psi(z_1),\ldots,\widehat\psi(z_n))
\]
is surjective for all sufficiently large \(R\). In particular, no finite orbit of evaluations in Weil's criterion remains invisible to the cofinal family.
\end{lemma}

\begin{proof}
The first assertion follows from compactness of the support. For the second, if the evaluation functionals were linearly dependent, there would be \(c_1,\ldots,c_n\), not all zero, such that
\[
 \int_{\mathbb R}\psi(x)
 \left(\sum_{j=1}^nc_je^{-iz_jx}\right)\,dx=0
 \qquad(\psi\in\mathcal D_R).
\]
The analytic function \(\sum_jc_je^{-iz_jx}\) would vanish on interval \((-R/2,R/2)\), hence throughout \(\mathbb C\). Independence of exponentials with distinct frequencies forces \(c_j=0\) for every \(j\), a contradiction. The functionals are independent; a linear map to \(\mathbb C^n\) with injective adjoint is surjective.
\end{proof}

\section*{Coupling of 2 routes and adversarial plateau test}

The bare product \(C_\angle\otimes I\) would require
\[
 2\|\psi\|^2\leq\|(I-T_\ell)\psi\|^2
\]
for every \(\psi\), an inequality that fails on smooth plateaux much longer than \(\ell\). Network \eqref{eq:amp-rh-red-prospectiva} makes no such identification: it preserves the Paley--Wiener, gamma, prime, and polar cross terms before projection. The following calculation exhibits that difference on a specific plateau.

Let \(\omega=e^{2\pi i/9}\) and, in the surviving fibre,
\[
 \kappa_0=u_{81}\otimes\chi_6,\qquad
 \kappa_1=v_{10}\otimes\chi_6,\qquad
 \chi_6(c)=3^{-1}\omega^{6c}.
\]
They are \(2\) orthogonal routes of the same role \(P_6\). Their boundary publication produces \(e_0,e_1\) with
\[
 \langle e_i,\Lambda_{\rm tw}^{\#}e_j\rangle=\delta_{ij}.
\]
Take
\[
 a=\frac1{16},\qquad b=2\log2,\qquad
 \psi_w=w^{-1/2}{\bf1}_{[-w/2,w/2]}\quad(w=a,b).
\]
These functions belong to the form closure of \(\mathcal D_R\), and their internal mollifications belong to \(\mathcal D_R\). If \(g_{ij}=\mathscr I_{\rm GW}(\psi_{w_i},\psi_{w_j})\), with \((w_0,w_1)=(a,b)\), the preceding expansion gives, for \(\lambda_n=2n+\tfrac12\),
\begin{align}
 g(w,w)
 ={}&\frac2w\sum_{n\geq0}
 \frac{1-e^{-\lambda_nw}}{\lambda_n^2}
 +c_\Gamma+\frac{32}{w}\sinh^2\!\frac w4 \nonumber\\
 &-2\sum_{k\log p<w}
  w_{p,k}\left(1-\frac{k\log p}{w}\right),
 \label{eq:amp-rh-meseta-diagonal}\\
 g_{01}
 ={}&\frac2{\sqrt{ab}}\sum_{n\geq0}
 \frac{e^{-\lambda_n(b-a)/2}(1-e^{-\lambda_na})}{\lambda_n^2}
 +c_\Gamma\sqrt{\frac ab} \nonumber\\
 &+2A_aA_b-\frac{\log2}{\sqrt2}\sqrt{\frac ab},
 \qquad A_w=\frac{4\sinh(w/4)}{\sqrt w}.
 \label{eq:amp-rh-meseta-cruce}
\end{align}
Bounding the geometric tails and the series of \(\lambda_n^{-2}\) yields
\begin{align}
 1.46610954095406&<g_{00}<1.46690960095857,\nonumber\\
 0.06966817455957&<g_{11}<0.06970424464086,\nonumber\\
 -0.008430670493497&<g_{01}<-0.008430670493494,
 \label{eq:amp-rh-meseta-cotas}
\end{align}
and, by interval multiplication,
\begin{equation}\label{eq:amp-rh-meseta-determinante}
 0.10207009921767<
 \det\begin{pmatrix}g_{00}&g_{01}\\g_{01}&g_{11}\end{pmatrix}
 <0.10217874948627.
\end{equation}
The complete Gram matrix of the plateau and the base state therefore has rank \(2\); a single scalar line cannot publish it.

Let us suppose
\[
 \sigma_b=g_{11}-\frac{g_{01}^2}{g_{00}}>0
\]
and define on \(\mathcal E_2=\operatorname{span}\{\psi_a,\psi_b\}\)
\begin{align}
 H^{(2)}_{m,R}(\alpha\psi_a+\beta\psi_b)
 ={}&
 e_{0,m}\left(
 \sqrt{g_{00}}\,\alpha
 +\frac{g_{01}}{\sqrt{g_{00}}}\,\beta\right)
 +e_{1,m}\sqrt{\sigma_b}\,\beta .
 \label{eq:amp-rh-lector-dos-rutas}
\end{align}
A Cholesky multiplication, using \(\langle e_i,\Lambda_m^\#e_j\rangle=\delta_{ij}\), proves
\begin{equation}\label{eq:amp-rh-meseta-energia-exacta}
 \langle H^{(2)}_{m,R}\psi,
         \Lambda_m^\#H^{(2)}_{m,R}\phi\rangle
 =\mathscr I_{\rm GW}(\psi,\phi)
 \qquad(\psi,\phi\in\mathcal E_2).
\end{equation}
The coefficients of \eqref{eq:amp-rh-lector-dos-rutas} come from the explicit series \eqref{eq:amp-rh-meseta-diagonal}--\eqref{eq:amp-rh-meseta-cruce}; no zeros are introduced, nor is the determinant's sign assumed: it is proved in \eqref{eq:amp-rh-meseta-determinante}. By continuity of the \(3\) channels, the identity and strict positivity of the determinant persist for sufficiently close smooth mollifications. This is the adversarial test that tensor \(C_\angle\otimes I\) fails and the indecomposable coupling passes.

\section*{Exact residual and regulator delimitation}

Define
\begin{equation}\label{eq:amp-rh-residual-finito}
 \mathcal R_{m,R,L}
 :=
 (\mathcal C^+_{m,R,L})^*\mathcal C^+_{m,R,L}
 -(\mathcal C^-_{m,R,L})^*\mathcal C^-_{m,R,L}
 -\mathscr I_{{\rm GW},L}.
\end{equation}
Lemma \ref{lem:amp-rh-expansion-guinand-weil} proves
\[
 \mathcal R_{m,R,L}=0
\]
term by term, not by definition. At the same time, \eqref{eq:amp-rh-publicaciones-construidas} gives, by computing norms, the port balance
\begin{equation}\label{eq:amp-rh-balance-residual-puerto}
 \boxed{
 \begin{aligned}
 \mathscr I_{{\rm GW},L}(\psi,\psi)
 -\|\mathsf L_{m,R,L}\psi\|^2
 ={}&2\operatorname{Re}
 \langle J_{-,m,R,L}\psi,\mathcal E_{m,R,L}\psi\rangle\\
 &+\|\mathcal E_{m,R,L}\psi\|^2 .
 \end{aligned}}
\end{equation}
Indeed, \(\|\mathsf L\psi\|^2
=\|J_+\psi\|^2-\|\mathscr C^{\rm src}J_+\psi\|^2\) and \(\mathscr C^{\rm src}J_+=J_-+\mathcal E\). Therefore, within this auxiliary chart, a proved vanishing of \(\mathcal E\) would yield the positive square
\begin{equation}\label{eq:amp-rh-cuadrado-regulado}
 \mathcal E_{m,R,L}=0
 \quad\Longrightarrow\quad
 \mathscr I_{{\rm GW},L}(\psi,\psi)
 =\|\mathsf L_{m,R,L}\psi\|^2\geq0.
\end{equation}
The implication does not authorize assuming its premise and is not part of the source chain of cofinal closure.

\begin{proposition}[Short-window obstruction]
\label{prop:amp-rh-obstruccion-ventana-corta}
For columns \eqref{eq:amp-rh-columna-mas-explicita}--\eqref{eq:amp-rh-columna-menos-explicita}, residual \(\mathcal E_{m,R,L}\) cannot vanish for every \(\psi\in\mathcal D_R\) and every \(L>0\).
\end{proposition}

\begin{proof}
Choose
\[
 0\ne\psi\in\mathcal D_R,\qquad
 \widehat\psi(i/2)=\widehat\psi(-i/2)=0.
\]
Such a function exists because two linear functionals are imposed on the infinite-dimensional space \(\mathcal D_R\). If \(2L<\log2\), the prime-power sum of \eqref{eq:amp-rh-gw-truncado-explicito} is empty and the polar channel vanishes. For \(h_{\psi,\psi}\), we have \(\check h_{\psi,\psi}(0)=\|\psi\|_2^2\), so that
\begin{equation}\label{eq:amp-rh-ventana-corta-expansion}
 \mathscr I_{{\rm GW},L}(\psi,\psi)
 =c_\Gamma\|\psi\|_2^2+
 \int_0^{2L}\vartheta_L(a)\mu_\Gamma(a)
       \|(I-T_a)\psi\|_2^2\,da .
\end{equation}
Since \(\|(I-T_a)\psi\|_2\leq a\|\psi'\|_2\) and \(\mu_\Gamma(a)=O(a^{-1})\) when \(a\downarrow0\), the integral of \eqref{eq:amp-rh-ventana-corta-expansion} is \(O_\psi(L^2)\). Because \(c_\Gamma<0\), for sufficiently small \(L\) we obtain
\[
 \mathscr I_{{\rm GW},L}(\psi,\psi)<0.
\]
If \(\mathcal E_{m,R,L}\psi=0\), implication \eqref{eq:amp-rh-cuadrado-regulado} would give the opposite sign. The residual therefore cannot vanish under the stated quantifier.
\end{proof}

Equality \(\mathcal R_{m,R,L}=0\) correctly identifies the difference of the two Gram matrices with the truncated functional; it does not turn it into a positive square. This regulated connector is therefore a \texttt{CONTROL\_NON\_GOVERNING}: it plays no part in the proof chain \eqref{eq:amp-rh-owner-binding}--\eqref{eq:amp-rh-owner-publicacion}. Proposition \ref{prop:amp-rh-obstruccion-ventana-corta} delimits only formulas \eqref{eq:amp-rh-columna-mas-explicita}--\eqref{eq:amp-rh-gw-truncado-explicito}; its result demotes neither cofinal positivity nor the \cref{thm:amp-rh-cierre-cofinal-propietario}.
